\begin{pastexam}[label={px:ntu-108-comm-dsab-32}]{臺大電信所 108 學年度（B）第 32 題}
（單選；每題 2.5 分，答錯不倒扣）

（第 31 題：Consider disjoint-set union/find algorithms with weighted union and path compression.
Assume the disjoint-set forest consists of $n$ items, what is the worst-case time complexity of
$n$ find/union operations? (choose the tightest bound)）

Following the previous question, the worst-case time complexity of one find operations is (choose
the tightest bound)

A.~$O(1)$ \qquad B.~$O(\log^* n)$ \qquad C.~$O(\log n)$ \qquad D.~$O(n)$ \qquad E.~$O(n\log n)$
\end{pastexam}

\begin{pastsoln}
答案是 C。

\textbf{上界}：依大小合併時，任何時候一棵有 $s$ 個節點的樹高度不超過 $\lfloor \lg s \rfloor + 1$（資結書
定理~13.2），路徑壓縮只會讓節點變淺，所以一次 \proc{Find} 最多走 $O(\lg n)$ 步。

\textbf{做得到}：每次都合併兩棵一樣大的樹，例如 $n = 8$ 時依序合併 $(1, 2)$、$(3, 4)$、$(5, 6)$、
$(7, 8)$、$(1, 3)$、$(5, 7)$、$(1, 5)$，這幾次合併都是對根做的，路徑壓縮沒有機會發揮，最後的樹高度
$\lg 8 + 1 = 4$（資結書第~13.3~節的例子）。 這時第一次 \proc{Find}$(8)$ 要經過 $4$ 個節點。 一般的
$n = 2^k$ 一樣得到深度 $\lg n$ 的節點，單次 \proc{Find} 是 $\Theta(\lg n)$。 路徑壓縮只對「之後」的
操作有幫助，壓縮那一次本身仍然要走完整條路徑。

A、B 是平攤的直覺（平均每次 $O(\alpha(n))$），不是單次最壞；D、E 太鬆。
\end{pastsoln}
